\documentclass[aspectratio=169]{beamer}
\input{header}
\coursecode{JKSSB}

\title{SMTP, POP \& IMAP}
\subtitle{Lesson 38 --- Mail Protocols and Threats}
\author{JKSSB Computer Awareness}
\date{\today}

\begin{document}
\frame{\titlepage}

\begin{frame}{Why Email Needs Protocols}
  \begin{itemize}
    \item Sending, receiving, and reading mail involve different mail
      servers and different programs.
    \item \key{Protocols} are fixed rules so any mail program can talk to
      any mail server.
    \item One protocol \key{sends} mail; two other protocols
      \key{receive} it in different ways.
    \item Secure transmission and threat awareness sit on top of these
      protocols.
  \end{itemize}
  \begin{block}{The one idea}
    SMTP sends mail out; POP3 and IMAP bring it in --- in two different ways.
  \end{block}
\end{frame}

\begin{frame}{SMTP: The Sending Protocol}
  \begin{itemize}
    \item \key{SMTP} stands for Simple Mail Transfer Protocol.
    \item It is the protocol used to \key{send} mail from the sender's
      program to the mail server, and between mail servers.
    \item The sender's program hands the message to its SMTP server; that
      server forwards it toward the recipient's mail server.
    \item SMTP does \key{not} deliver mail to the recipient's inbox for
      reading --- receiving is the job of POP3 or IMAP.
  \end{itemize}
  \begin{alertblock}{Direction to remember}
    SMTP is for \key{sending}: out of the sender, toward the
    recipient's server.
  \end{alertblock}
\end{frame}

\begin{frame}{POP3: Download to One Machine}
  \begin{itemize}
    \item \key{POP3} stands for Post Office Protocol version 3.
    \item It is a \key{receiving} protocol: it pulls mail from the mail
      server down to one local machine.
    \item After downloading, the mail is typically \key{removed from the
      server}, so it lives on that one computer.
    \item It suits a single computer where mail is read offline in one
      place.
  \end{itemize}
  \begin{block}{Keep the pair straight}
    POP3 \key{downloads} mail to the local machine; it does not keep it
    synced on the server.
  \end{block}
\end{frame}

\begin{frame}{IMAP: Sync Across Devices}
  \begin{itemize}
    \item \key{IMAP} stands for Internet Message Access Protocol.
    \item It is also a \key{receiving} protocol, but it works
      differently from POP3.
    \item Mail \key{stays on the server}; the program shows a synced
      view of folders, read marks, and deletions.
    \item Reading, deleting, or moving a message on the phone shows the
      same change on the laptop, because every device mirrors the
      server.
  \end{itemize}
  \begin{block}{Office desktop example}
    Checking office mail on the phone and later on the desktop shows the
    same inbox --- that is IMAP keeping every device in step with the
    server.
  \end{block}
\end{frame}

\begin{frame}{POP3 vs IMAP}
  \begin{center}
    \begin{tabular}{lll}
      \toprule
      \textbf{Point} & \textbf{POP3} & \textbf{IMAP} \\
      \midrule
      Job & Receives mail & Receives mail \\
      Method & Downloads to one machine & Syncs a view of the server \\
      Server copy & Usually removed & Kept on the server \\
      Many devices & Poor fit & Designed for it \\
      \bottomrule
    \end{tabular}
  \end{center}
  \vspace{0.35cm}
  \begin{alertblock}{Trap alert}
    POP3 \key{downloads} to local storage; IMAP \key{syncs} with the
    server. Do not swap them. (TRAP-17)
  \end{alertblock}
\end{frame}

\begin{frame}{Mail-Server Flow, End to End}
  \begin{itemize}
    \item The sender writes a message and the mail program sends it by
      \key{SMTP} to the sender's mail server.
    \item The sender's server forwards it by \key{SMTP} to the
      recipient's mail server, where it waits in the mailbox.
    \item The recipient's program collects it with \key{POP3} (download)
      or \key{IMAP} (sync).
    \item Sending is always SMTP; the receiving side is POP3 or IMAP.
  \end{itemize}
  \begin{block}{The flow in one line}
    Sender $\rightarrow$ SMTP $\rightarrow$ recipient server
    $\rightarrow$ POP3 or IMAP $\rightarrow$ reader.
  \end{block}
\end{frame}

\begin{frame}{To, Cc, and BCC}
  \begin{itemize}
    \item \key{To} names the main recipients; \key{Cc} (carbon copy)
      names others kept in the loop.
    \item \key{BCC} stands for blind carbon copy: its recipients are
      \key{hidden from each other}.
    \item Everyone sees the To and Cc lists; nobody on the To or Cc
      list sees who is on BCC.
    \item The \key{sender} always knows all recipients; BCC hides
      recipients from each other, not from the sender.
  \end{itemize}
  \begin{alertblock}{Trap alert}
    BCC hides recipients \key{from each other}. (TRAP-16)
  \end{alertblock}
\end{frame}

\begin{frame}{Filters, Rules, and Signatures}
  \begin{itemize}
    \item \key{Filters (rules)} sort incoming mail automatically ---
      for example, moving newsletters to one folder or flagging
      important senders.
    \item A filter acts on the \key{headers and content} of the message,
      such as sender, subject, or words in the body.
    \item A mail \key{signature} is a fixed block (name, post, phone)
      added at the end of outgoing messages.
    \item A signature identifies the sender; it does \key{not} encrypt
      or protect the message.
  \end{itemize}
  \begin{block}{Keep the pair straight}
    Filters organise incoming mail; a signature labels outgoing mail.
    Neither one encrypts anything.
  \end{block}
\end{frame}

\begin{frame}{Digital Signature vs Encryption}
  \begin{itemize}
    \item A \key{digital signature} proves who sent the message and that
      it was not changed --- authenticity and integrity.
    \item \key{Encryption} scrambles the message so only the intended
      recipient can read it --- confidentiality.
    \item A signed message can still be \key{read} by anyone; a signed
      message is not automatically secret.
    \item An encrypted message can still be \key{forwarded}; encryption
      hides content, it does not prove the sender.
  \end{itemize}
  \begin{alertblock}{Trap alert}
    A digital signature does \key{not} encrypt the body. Signature =
    authenticity; encryption = secrecy. (TRAP-15)
  \end{alertblock}
\end{frame}

\begin{frame}{Secure Email: SSL and TLS}
  \begin{itemize}
    \item \key{SSL} stands for Secure Sockets Layer; \key{TLS} stands
      for Transport Layer Security.
    \item Together \key{SSL/TLS} provide the secure channel over which
      mail is transmitted.
    \item Without this protection, mail travels in a form others on the
      path could read.
    \item Exam answer: the secure email transmission protocol is
      \key{SSL/TLS}. (PYQ-36)
  \end{itemize}
  \begin{block}{Exam recall}
    Secure transmission of email = SSL/TLS.
  \end{block}
\end{frame}

\begin{frame}{Spam}
  \begin{itemize}
    \item \key{Spam} is bulk unwanted mail, usually advertising or junk
      sent to many addresses at once.
    \item It clutters the inbox and can carry fraud or malware, but the
      defining mark is \key{bulk and unwanted}.
    \item Spam filters move suspected junk to a separate folder using
      sender reputation and content rules.
    \item \muted{Deleting spam without opening links or attachments is
      the safe habit.}
  \end{itemize}
  \begin{exampleblock}{Office desktop example}
    Dozens of identical ``prize winner'' messages arriving overnight go
    to the junk folder --- that is spam filtering at work.
  \end{exampleblock}
\end{frame}

\begin{frame}{Phishing}
  \begin{itemize}
    \item \key{Phishing} is a fake message that pretends to be a trusted
      sender --- a bank, an office, or a colleague --- to steal
      passwords or money.
    \item It creates \key{urgency}: ``verify now'', ``account blocked'',
      ``claim within 24 hours''.
    \item It asks the reader to click a link or open an attachment and
      hand over secrets.
    \item Check the real sender address and never enter passwords from a
      mail link.
  \end{itemize}
  \begin{alertblock}{Takeaway}
    Phishing = impersonation + urgency + a link or attachment that
    steals secrets.
  \end{alertblock}
\end{frame}

\begin{frame}{Spoofing, Malicious Links, and Attachments}
  \begin{itemize}
    \item \key{Spoofing} forges the sender address so a message looks as
      if it came from someone trusted.
    \item A \key{malicious link} leads to a fake login page or a
      malware download; hovering shows an address that does not match
      the claimed sender.
    \item A \key{malicious attachment} carries malware behind an
      innocent-looking file name.
    \item Rule: verify the sender by another channel before clicking or
      opening anything unexpected.
  \end{itemize}
  \begin{block}{Keep the pair straight}
    Phishing is the \key{fraud attempt}; spoofing is the \key{forged
    sender} trick it often uses.
  \end{block}
\end{frame}

\begin{frame}{Trap Alert: Mail Facts to Keep Straight}
  \begin{itemize}
    \item SMTP \key{sends}; POP3 and IMAP \key{receive}.
    \item POP3 \key{downloads} to one machine; IMAP \key{syncs} with
      the server. (TRAP-17)
    \item A digital signature proves the sender; it does \key{not}
      encrypt the message. (TRAP-15)
    \item BCC hides recipients \key{from each other}, not from the
      sender.
    \item Secure email transmission = \key{SSL/TLS}. (PYQ-36)
  \end{itemize}
\end{frame}

\begin{frame}{Lesson 38: Recall}
  \begin{itemize}
    \item SMTP (Simple Mail Transfer Protocol) sends mail; POP3 (Post
      Office Protocol version 3) downloads it to one machine; IMAP
      (Internet Message Access Protocol) syncs it across devices.
    \item Mail flow: sender $\rightarrow$ SMTP $\rightarrow$
      recipient server $\rightarrow$ POP3 or IMAP $\rightarrow$
      reader.
    \item BCC hides recipients from each other; filters and rules sort
      incoming mail; a signature labels outgoing mail.
    \item Digital signature = authenticity and integrity; encryption =
      confidentiality; secure transmission = SSL/TLS.
    \item Spam is bulk unwanted mail; phishing impersonates trust to
      steal secrets; spoofing forges the sender; malicious links and
      attachments deliver the fraud or malware.
  \end{itemize}
\end{frame}
\end{document}
